\section*{Appendix C: Service Designer’s Quick-Start}

To instantiate a new Service Node within an existing Society, a Designer must follow the "Instantiation Protocol":

\begin{enumerate}
    \item \textbf{Identify Latent Demand:} Use the local \textit{Prospecting Engine} to verify that a community need is currently unmet (e.g., a gap in local renewable energy maintenance).
    \item \textbf{Draft the WBS:} In the Modeler, drag and drop standardized \textit{Task Blocks} from the relevant BoK. 
    \item \textbf{Define Environment Properties:} Set the "Trigger Events" (e.g., \textit{If Grid\_Voltage < 210V, trigger Maintenance\_Task}).
    \item \textbf{Determine Resource Mapping:} Categorize every required input as either \textit{Internal} (sourced from the Mesh) or \textit{External} (requiring a Common Fund lease).
    \item \textbf{Ratification:} Present the model to the local Society. Stake the required \textit{Reputation Tokens} to commit the WBS to the local ledger.
\end{enumerate}



\noindent \textit{Note: Designers are encouraged to "Fork" existing high-efficiency Node models from the Global BoK to accelerate community bootstrapping.}
